% !PRINTED 73
\weberpage{73}

Then the $\gamma$ must satisfy the condition:

\weq{1}{C_{-1}^{\gamma_{-1}} C_0^{\gamma_0} C_1^{\gamma_1} \dots C_\mu^{\gamma_\mu} \equiv 1 \, (\mathrm{mod}\ m').}

This condition demands that all indices, with the exception of the one corresponding
to the prime $q$, which we denote by $\gamma$, must be zero. With respect to
$\gamma$, however, a distinction is to be made according as $q$ still divides $m'$
or not, i.\,e. whether $q$ is a multiple or only a simple factor of $m$. If $q$ is
no longer contained in $m'$, then the congruence (1) contains no restriction at all
for $\gamma$, and $\gamma$ can take every value with respect to the modulus $c$,
i.\,e. every value $0, 1, \dots q - 2$.

But if $q$ is still contained in $m'$, and hence $m$ is divisible by $q^\varkappa$, and $\varkappa > 1$,
then the condition (1) demands:
\[
C^\gamma \equiv 1 \, (\mathrm{mod}\ q^{\varkappa-1}),
\]
hence
\weq{2}{\gamma \equiv 0 \, \left(\mathrm{mod}\ \dfrac{c}{q}\right),}
and $\gamma$ can therefore receive each of the values
\[
0, \dfrac{c}{q}, \dfrac{2c}{q}, \dots \dfrac{(q-1)c}{q}
\]
and no others. In the first case the degree of the group $Q$ is equal to $q - 1$, in the
second equal to $q$.

The indices $\delta$ of a number from the group $P$ are obtained, by \S.\,12, 7.
and \S.\,16, (7), from the condition that for all admissible $\gamma$
\weq{3}{\varepsilon_{-1}^{\gamma_{-1}\delta_{-1}} \varepsilon_0^{\gamma_0\delta_0} \varepsilon_1^{\gamma_1\delta_1} \dots \varepsilon_\mu^{\gamma_\mu\delta_\mu} = 1}
shall hold. By what has just been proved about the indices $\gamma$, however,
(3) demands only the one thing, that the index $\delta$ corresponding to the
prime $q$ shall be divisible by $q$ or by $q - 1$, according as $q$ divides $m$
several times or only once. We therefore emphasise the theorem:

7. \gesperrt{The group $P$ reciprocal to $Q$ is characterised by the fact that
the index $\delta$ corresponding to $q$ of all its numbers is divisible by $q$
or by $q - 1$, according as $q$ divides $m$ several times or only once.}

And from this:

8. \gesperrt{A non-primary divisor $\mathfrak{A}$ of $\mathfrak{N}$ is
characterised by the fact that in the reciprocal group $\mathfrak{B}$}
%CONT>>
